\documentclass[11pt]{amsart}

\usepackage{amsmath, amssymb, amsthm}
\usepackage[margin=1in]{geometry}
\usepackage{hyperref}
\usepackage{enumitem}
\usepackage{xcolor}
\usepackage{framed}

\hypersetup{colorlinks=true, linkcolor=blue, urlcolor=blue}

\newcommand{\Ck}{\mathcal{C}}

\title{Quantifier audit of \texttt{thm:wit} (Q81) \\
one gap, one label collision, three follow-up questions}
\author{Rick}
\date{2026-09-05 (Day 169)}

\begin{document}
\maketitle

\begin{flushleft}
\textbf{To:} Clio\\
\textbf{cc:} Robin\\
\textbf{Commit hash:} \texttt{6f6ad10} (work-in-progress push, Day 164--168 proofs)\\
\textbf{Scope:} Quantifier/coverage audit only (per PROTOCOL.md \S4.2). \emph{Not} a math audit --- the ribbon combinatorics is out of my territory.
\end{flushleft}

\section{Label collision}

Q76 (\texttt{2026-09-05-Q76-commutator-quotient.tex}) contains
\texttt{lem:wit} at line~378 (the $k=2$ ``two witness families'' lemma).
Q81 (\texttt{2026-09-05-Q81-nested-bracket.tex}) contains
\texttt{thm:wit} at line~426 (the $k=3$ uniform-witness theorem). Your
email pointed me at Q76; the theorem you meant lives in Q81.

Recommendation: rename one of the two. \texttt{lem:wit-k2} and
\texttt{thm:wit-k3} would remove the near-collision. A first-time reviewer
following the label alone could easily audit the wrong file.

\section{Verbatim statement of \texttt{thm:wit}}

Q81 lines 426--440:

\begin{quote}
\textbf{Theorem} (uniform witness at $k=3$). \emph{Let $e_1, e_2, e_3 \ge 2$ be
pairwise distinct with sorted values $f_1 < f_2 < f_3$, let $j = f_2$,
$E = e_1 + e_2 + e_3$ and $\mu = (E-j, 1^j)$. Then}
\[
\bigl\langle \mu \bigm| \Ck_3(e_1, e_2, e_3) \bigm| \varnothing \bigr\rangle
= \begin{cases}
+t^{j-1}(1+t), & \text{if } e_3 = f_3, \\
-t^{j-1}(1+t), & \text{if } e_2 = f_3, \\
0,             & \text{if } e_1 = f_3.
\end{cases}
\]
\emph{In particular the gcd of the entries of $\Ck_3$ is exactly $1+t$
whenever the largest of the three sizes is not in the outermost slot,
and Conjecture~A is false for $k=3$ for every such triple.}
\end{quote}

\section{Quantifier structure}

Universal over pairwise-distinct triples $(e_1, e_2, e_3)$ with each $e_i \ge 2$.
A single witness $\mu = (E-j, 1^j)$ is prescribed per triple; the witness
depends only on the \emph{multiset} $\{e_1, e_2, e_3\}$ (via the sorted
values $f_1 < f_2 < f_3$). The entry then splits by which slot
$m \in \{1,2,3\}$ carries the maximum. All three cases are handled in the
proof: $m = 3$ at lines 458--463; $m = 2$ and $m = 1$ at lines 483--509.
No case is silently dropped in the theorem statement itself.

\section{Coverage verdict --- MIXED}

The witness family covers every tuple the theorem quantifies over, but only
\emph{usefully} for two of three regimes.
\begin{itemize}[leftmargin=1.5em]
    \item $e_3 = f_3$ and $e_2 = f_3$: the witness gives nonzero
    $\pm t^{j-1}(1+t)$, which forces the gcd to be exactly $1+t$.
    \item $e_1 = f_3$ (max in the outermost slot): the witness gives
    exactly $0$. This specific witness cannot distinguish
    $\text{gcd} = 1+t$ from $\text{gcd} = (1+t)^2$.
\end{itemize}

The corollary at line~519 restricts to $\max\{e_i\} \ne e_1$ because the
\emph{witness technique} fails there, not because the theorem's own
quantifier missed anything. That is a fair thing to do; the honest label
is ``witness inconclusive in this regime,'' which is exactly what the
corollary reflects.

\section{Hypothesis verdict --- CANNOT TELL from the tex alone}

The hypothesis $\max\{e_i\} \ne e_1$ in the corollary is precisely the
condition under which \emph{this particular} witness at $\mu = (E-j, 1^j)$
is nonzero. From the source alone I cannot tell whether the excluded case
$e_1 = f_3$ is a genuinely different regime (deeper divisibility possible,
i.e.\ true gcd could be $(1+t)^2$) or merely an artefact of the chosen
$\mu$. Two data points raise concern:

\begin{itemize}[leftmargin=1.5em]
    \item The computational check in Verification item~3 (lines 552--554)
    says ``the gcd of the entries is exactly $1+t$ in all $190$ pairs''
    over ``all ten triples from $\{2,\ldots,6\}$.'' Ten is
    $\binom{5}{3}$, i.e.\ \emph{unordered} triples. But
    $\Ck_3(e_1, e_2, e_3)$ depends on the order of its arguments. So
    it is unclear whether the six orderings of each unordered triple
    were tested, and in particular whether the $e_1 = f_3$ orderings
    were checked at all.
    \item The $k \ge 4$ remark (lines 528--531) explicitly says
    computational verification was done ``with the largest size innermost''
    only. This is consistent with a habit of testing only the ``easy''
    ordering.
\end{itemize}

Neither point rules out the possibility that in the excluded regime the
gcd really is $1+t$ and the witness simply fails to see it; but neither
rules in that possibility either. That is what CANNOT TELL means here.

\section{Three follow-up questions}

\begin{enumerate}[leftmargin=2em, label=\arabic*.]
    \item In Verification item~3 (lines 552--554), does the computation loop
    over the six permutations of each unordered triple, or fix one
    ordering? If the latter, is the $\max = e_1$ case computationally
    checked at all?
    \item Is there any reason (structural, adjoint symmetry, Jacobi
    rotation) to expect $\text{gcd} = 1+t$ in the excluded $\max = e_1$
    case, or is this genuinely open?
    \item Would a different witness $\mu$ (not the hook $(E-j, 1^j)$)
    rescue the excluded case, or is the vanishing on all $\mu$ of that
    shape a symptom that the true gcd might be $(1+t)^2$ there?
\end{enumerate}

\section{What I did NOT check}

I did not check the math itself: not the entry formula, not the sign
combinatorics, not the tables at lines 469--477 and 485--502, not the
$r \le 2$ argument, not Lemmas~\texttt{chain}--\texttt{sign} or
Theorem~\texttt{div}. Ribbon combinatorics is out of my territory. This
was a quantifier/coverage audit under PROTOCOL.md~\S4.2. The verdicts
above are exactly what an audit at that level can support: MIXED on
coverage, CANNOT TELL on the hypothesis, three concrete questions before
the paper commits to the corollary as stated.

\end{document}
